% !TEX root = ../../main.tex

\section{本章小结}

本章在需求分析的基础上，对 HearBridge 系统进行了总体设计。
首先，明确了系统设计目标与原则，提出前后端分离、多端协同、识别服务独立、资源与模型可维护以及有限场景可控等设计原则；
其次，设计了系统总体架构，说明了 HarmonyOS 移动端、Vue Web 管理端、Spring Boot 服务端、Python 手势识别服务、DeepSeek 语义修正服务以及 MySQL、Redis、MinIO 等支撑组件之间的关系；
随后，对系统功能模块进行了划分，分别说明了移动端、后台管理端、服务端和手势识别服务的模块职责；
接着，完成了核心业务 E-R 图、主要数据表设计和持久层设计。

通过本章设计，可以明确 HearBridge 系统的整体结构、功能边界和核心数据关系。
下一章将在系统设计的基础上，进一步说明各端核心功能的具体实现过程。
